problem:  
Let \(\mathfrak{C}_{\mathrm{sec}>0}\subset \mathrm{Sym}^2(\Lambda^2\mathbb{R}^n)\) denote the algebraic cone of curvature tensors with **positive sectional curvature** in dimension \(n\ge 4\).  The Ricci flow curvature operator evolution equation induces a nonlinear vector field  
\[
Q(R)=R^2 + R^{\#},
\]
where \(R^{\#}\) is the quadratic term arising from the curvature tensor bracket.  

It is known that some cones (e.g., positive curvature operator, two‑positive curvature operator, positive isotropic curvature) are invariant under this evolution, but the invariance of \(\mathfrak{C}_{\mathrm{sec}>0}\) is not established and is believed delicate.  

**Open problem:**  
Determine whether there exists a *minimal algebraic curvature cone* \(\mathfrak{C}_{\min}\) satisfying
\[
\mathfrak{C}_{\mathrm{sec}>0} \subsetneq \mathfrak{C}_{\min} \subset \{\text{curvature operators with }\mathrm{Ric}>0\}
\]
such that \(\mathfrak{C}_{\min}\) is **invariant under the Ricci flow ODE** \(\dot R = Q(R)\).  

Equivalently: can one describe, at the curvature‑tensor level, the smallest Ricci‑flow–invariant cone containing all positively sectional curved tensors?  
If such a cone exists, does it coincide with one of the known Wilking invariant cones, or is it genuinely new?  

A positive answer would identify a new algebraic layer controlling the passage from **local sectional curvature geometry** to **Ricci flow invariance**, while a negative one would indicate that no finite curvature pinching condition bridges the gap between positive sectional curvature and known invariant cones.

Why is it a "good" problem:  
This reframing converts the classical, global curvature‑preservation question into a **foundational algebraic‑geometric analysis problem** within the Ricci flow’s curvature ODE. It asks for a structural synthesis between **algebraic curvature inequalities** and **analytic invariance under nonlinear evolution**—a question that naturally extends Böhm–Wilking’s curvature cone framework.  
The problem is simple to state but difficult: it demands a complete understanding of the interplay between the quadratic term \(R^{\#}\) and the nonlinear geometry of cone boundaries. Resolving it would clarify the precise algebraic threshold where positive sectional curvature loses analytic control under Ricci flow, thus tightly linking **invariant theory, nonlinear PDE evolution, and Riemannian geometry**. This bridges core aspects of geometric analysis and differential geometry in a truly deep and generative way.